\FloatBarrier
\section{Moller cone tip}
\label{sec:cone-tip}

The eleventh beamline mesh study renames \texttt{cone-tip} to \texttt{moller\_cone\_tip} in all four source,
runtime, and GEMC3 CAD configurations: runs 5874, 11620, 18304, and 20000. The \texttt{beamline\_W} material,
passive 180-degree y rotation, and 1273.27 mm z translation remain unchanged. SQLite banks continue to use
\texttt{default}, indexed by run number. The original is archived; the three installed meshes have identical
bytes.

The original is a clean closed manifold solid with one component and genus two. Protected quadric decimation
reduces 13,730 triangles to 5,000 (63.5834\% fewer), and 686,584 bytes to 250,084 (63.5756\% fewer). All 1,936
faces adjoining edges above five degrees retain their exact triangle coordinates. The original bounds, central
passage, and other opening are retained; no hole filling or reconstruction is applied. The result has no
boundary edges, non-manifold elements, detected self-intersections, or zero-area triangles. The current
builder reproduces the saved STL byte for byte.

Volume changes from 196511.847100 to 196510.892015 mm$^3$, a change of -0.000486019\%. Bidirectional vertex,
edge, and surface sampling gives a maximum deviation of 0.032396 mm, below the 0.05 mm guard. This is a
sampled maximum, not a proven global error bound. More aggressive protected trials distort the geometry and
are rejected. Table~\ref{tab:cone-tip-mesh} gives the retained mesh properties.

\begin{table}[htbp]
\centering
\caption{Moller cone tip before and after protected decimation.}
\label{tab:cone-tip-mesh}
\begin{tabular}{lrr}
\toprule
Quantity & Original & Optimized \\
\midrule
Triangles & 13,730 & 5,000 \\
Bytes & 686,584 & 250,084 \\
Boundary edges & 0 & 0 \\
Components / genus & 1 / 2 & 1 / 2 \\
Volume (mm$^3$) & 196511.847100 & 196510.892015 \\
Detected intersecting faces & 0 & 0 \\
\bottomrule
\end{tabular}
\end{table}

The isolated benchmark selects the run-5874 CAD row with \texttt{default}; the tip is absent from run 11. The
benchmark runner now accepts \texttt{--runno} so it can select such components without remapping their run
keys. Both fixtures use matched 5 GeV electron, proton, positive-pion, and neutron primaries from the target,
uniform in solid angle over 2.5--4.8 degrees and full azimuth. Physics is FTFP\_BERT with no magnetic field or
optical processes, one thread, default production cuts, and no GUI or event output. Five alternating-order
repetitions use 200 events per particle and seeds 12345--12349. The fixed GEMC3 executable and its SHA-256
match the recent beamline studies; all primary hashes match within seed pairs.

Table~\ref{tab:cone-tip-tracking} reports median Geant4 beamOn wall times, excluding initialization. Ratios
are original/optimized; values above one mean faster tracking. All five pairs are faster for each particle,
and all 52 isolated benchmark processes complete without Geant4 exceptions. The sample measures this one mesh;
it does not establish complete CLAS12 speedup or detector-response equivalence.

\begin{table}[htbp]
\centering
\caption{Isolated Moller cone tip tracking; five repetitions of 200 events per particle.}
\label{tab:cone-tip-tracking}
\setlength{\tabcolsep}{6pt}
\begin{tabular}{lrrrr}
\toprule
Particle & Original (s) & Optimized (s) & Speedup & Paired ratio range \\
\midrule
Electron & 17.63610 & 10.63150 & 1.659 & 1.658--1.665 \\
Proton & 1.84329 & 1.62666 & 1.133 & 1.133--1.211 \\
Positive pion & 2.40739 & 2.04576 & 1.177 & 1.173--1.236 \\
Neutron & 1.62514 & 1.40266 & 1.159 & 1.159--1.207 \\
\bottomrule
\end{tabular}
\end{table}

Median loading process wall time is nearly unchanged, 0.843174 to 0.841411 s, while loading CPU time falls
from 0.80 to 0.76 s. Engine initialization wall time changes from 0.758880 to 0.727224 s. Startup peak RSS
falls from 133.125 to 122.813 MiB (7.75\%). Process CPU time, total wall time, memory, seed-paired ratios, and
raw diagnostics are retained separately from the tracking measurements.

The main-cone/tip assembly already overlaps before decimation. Independent surface collisions are present with
both engineering originals as well as with the current main cone. A separate 10,000-point Geant4 assembly
check reports five GeomVol1002 exception blocks both before and after the tip change; sampled penetrations
reach about 0.372 mm. This optimization preserves the engineering interfaces and does not repair that assembly
overlap. Collision counts and vertex-containment counts depend on tessellation and must not be interpreted as
overlap volumes or as evidence of improved clearance.

All four complete affected configurations construct 39 Geant4 volumes including the world. Their two
GeomSolids1001 warnings remain, and sampled overlap blocks change from 20 to 21 per construction. The renamed
volume changes construction order and the resulting sampled reports; contacts among unchanged parts remain.
The clean isolated benchmark therefore does not certify the complete assemblies. All 17 beamline
configurations pass volume-count construction checks. All 15 native-geometry, 15 material, and 14 CAD semantic
comparisons pass; direct database checks allow only the four names and STL paths to change. Every other mesh
retains its prior hash. Figure~\ref{fig:cone-tip} uses matched camera and bounds.

\begin{figure}[p]
\centering
\includegraphics[width=\linewidth]{figures/moller_cone_tip_stl_comparison.pdf}
\caption{Original and optimized Moller cone tip meshes with identical bounds and camera. Openings and
sharp features are retained while redundant smooth-surface triangles are removed.}
\label{fig:cone-tip}
\end{figure}
